\section{Scope and further questions}\label{sec:outlook}
The proved transfer has two distinct scales. A fixed number of free vertices contributes one factor $\Phi(\ell)^{-1}$ per vertex. At $f\asymp\sqrt M$, the cumulative mean shift survives the Gaussian normalization, and the positive quadratic weight produces a collective factor. The bounded offset of the degree cap contributes a linear correction on that same scale. Both are present in \eqref{eq:criticalexplicit}.

Several extensions remain meaningful and require new estimates.
\begin{enumerate}
\item \textbf{An effective relative remainder.} The present proof uses a uniform $o(1)$ dense enumeration theorem and qualitative limit arguments. An explicit remainder with computable constants would turn the inverse brackets into finite-input certificates. It would need quantitative control of the enumeration input, relative graph tails, tilted moments, and parity error together.
\item \textbf{Higher count coefficients.} The computable expansion of the saddle mass is only one ingredient. A next term for $A_n$ requires a correspondingly precise mixed degree-sequence expansion and all associated lattice and moment corrections. Expanding the known leading model alone does not produce a justified coefficient for the actual count.
\item \textbf{More than order $\sqrt M$ free vertices.} If $f/\sqrt M\to\infty$, the normalized mean shift $a_M$ is unbounded. The compact Gaussian argument and bounded-shift integrability used here no longer prove \eqref{eq:criticalexact}. A moving saddle adapted to the free fraction, together with moderate-deviation estimates, is a natural route; no extrapolated formula is claimed.
\item \textbf{Nonuniform caps and marked free degrees.} The graph concentration switching already permits unequal caps. Extending the summation to bounded heterogeneous offsets or generating functions for the free degrees would require choosing and analyzing the corresponding tilt. It may reveal which aggregate offset statistics govern the first collective correction.
\item \textbf{Historical identification.} The recovered historical table and direct loopy-graph predecessor narrow the attribution question, but they do not replace inspection of Anderson's cited chapter or the underlying unpublished Sillke--Flammenkamp work. Further access could clarify the history without changing the proof given here.
\end{enumerate}

The analytic conclusions of this report are therefore the leading relative matrix equivalent, the fixed and bounded-critical free-vertex transfer, and the inverse statements justified by their precision. The exact finite computations support the model and implementation, while the source record identifies the imported ingredients. Neither numerical agreement nor a bounded literature search is a priority certificate.
